\chapter*{前言}\addcontentsline{toc}{chapter}{前言}
线性代数的计算可以归结为少数操作，理解却依赖对同一对象的多种观察。一个矩阵既是数字排成的表，也是空间的线性变换；方程组既能用消元求解，也是在寻找若干约束的交集；二次型既能配方，也能从椭圆的主轴看出正定性与极值。

本册以六章考研核心为主线：行列式、矩阵、向量组、线性方程组、特征值与相似、二次型。第七章把计算提升为线性空间与线性变换；第八章集中处理秩不等式、矩阵多项式、Jordan结构、矩阵分解与范数。拓展内容设有明确标记，便于按基础与目标取舍。

贯穿全书的第一条线索是\textbf{维数}。秩在数独立方向，核在找被消去的方向，基础解系在列出剩余自由度。遇到复杂矩阵时，先判断哪些方向新增、哪些方向丢失，往往能把长计算压缩成一个维数公式。

第二条线索是\textbf{换一种观察坐标}。选好基，线性变换可能变成几个独立缩放；选好正交基，二次型的交叉项可能消失；数据中心化以后，最小二乘的联立方程也能简化。换基本身不是目标，目标是找到使结构显露的坐标。

第三条线索是\textbf{先写条件，再用结论}。可逆、对称、正定、满列秩和数域分裂都不是可有可无的修饰。忘记条件，常会把半正定判据写错，把重根误判成不可对角化，或在含参数的消元中丢掉整个分支。本册的易错提示着重解释这些错误为何发生，以及怎样用反例与代回检查发现它们。

全书保留128个完整例题与72道分层习题，并增加17幅数学图解。每幅图都对应一个概念、推导或解题动作：剪切说明行列式倍加不变，投影三角形说明最近性，核增长说明Jordan链，单位圆变椭圆说明奇异值。图形用于发现结构；精确数值与一般结论仍由定义、公式和证明确认。

公式按“条件、变形、理由、结论”展开：消元保留关键矩阵变化，参数题单列退化分支，换基交代坐标方向，多解题分别算到底。习题附录逐题给出完整解答；Jordan存在性、矩阵分解与伪逆等选学内容也配有推导，公式速查附录提供证明位置与补充计算。

“巧解”并不意味着绕过基本功。观察到低秩就写外积，看到相同背景就寻找秩一更新，给出多项式关系就先降次，这些方法都来自对条件的熟悉。学习多解时应比较适用范围、计算成本和能否推广，而不只是记录一个更短的答案。

希望读者每次计算结束后再做一次复盘：本题真正用了哪个条件？图形能否解释结果？若改变一个条件，方法是否仍成立？能回答这几问，方法才会成为可以迁移的能力。
